% sec:kerr-bl — Boyer-Lindquist Coordinates
\section{Boyer-Lindquist Coordinates}
\label{sec:kerr-bl}

Every astrophysical black hole for which we have observational data
carries angular momentum.  The unique stationary, axisymmetric vacuum
solution to the Einstein field equations is the Kerr metric, governed
by two parameters: the mass~$M$ and the specific angular momentum
$a = J/M$.  Cosmic censorship requires $\abs{a} \leq M$; the upper
bound $\abs{a} = M$ marks the extremal limit.
\coderef{Spacetime}{KerrKS}
\histnote{Roy Kerr found this solution in 1963 by studying
algebraically special spacetimes \cite{Kerr1963}.  Boyer and
Lindquist recast it into the coordinate system used here
\cite{BoyerLindquist1967}.}

\paragraph{Structural functions.}

The Boyer--Lindquist form of the Kerr metric is built from two
auxiliary functions that encode how rotation deforms the
Schwarzschild geometry.  The first measures the oblate-spheroid shape
of surfaces of constant~$r$:
%
\begin{equation}
  \Sigma = r^2 + a^2 \cos^2\!\theta \,.
  \label{eq:kerr-Sigma}
\end{equation}
%
Setting $a = 0$ gives $\Sigma = r^2$, recovering spherical symmetry.
For $a \neq 0$, $\Sigma$ vanishes only at $r = 0$, $\theta = \pi/2$
--- the location of the Kerr ring singularity, where the Kretschmann
invariant
$R_{\alpha\beta\gamma\delta}R^{\alpha\beta\gamma\delta}$ diverges
(\cref{sec:dg-curvature}).  The second function governs the radial
structure:
%
\begin{equation}
  \Delta = r^2 - 2Mr + a^2 \,.
  \label{eq:kerr-Delta}
\end{equation}
%
The roots $r_\pm = M \pm \sqrt{M^2 - a^2}$ locate the outer (event)
and inner (Cauchy) horizons; both exist precisely when
$\abs{a} \leq M$.

\paragraph{The line element.}

With $\Sigma$ and $\Delta$ in hand, we also define an auxiliary
function that appears in the azimuthal part of the metric:
%
\begin{equation}
  A = (r^2 + a^2)^2 - a^2 \Delta \sin^2\!\theta \,.
  \label{eq:kerr-A}
\end{equation}
%
The quantity $A/\Sigma$ governs the circumferential radius in the
azimuthal direction; its excess over $(r^2 + a^2)^2$ measures
how strongly rotation inflates the equatorial geometry.  The full
Boyer--Lindquist line element is
%
\begin{multline}
  \d s^2 = -\!\left(1 - \frac{2Mr}{\Sigma}\right)\d t^2
           - \frac{4Mar\sin^2\!\theta}{\Sigma}\,\d t\,\d\varphi
           + \frac{\Sigma}{\Delta}\,\d r^2 \\
           + \Sigma\,\d\theta^2
           + \frac{A\sin^2\!\theta}{\Sigma}\,\d\varphi^2 \,.
  \label{eq:kerr-bl-metric}
\end{multline}
%
The off-diagonal $\d t\,\d\varphi$ term --- absent in the
Schwarzschild metric (\cref{eq:schwarzschild-metric}) --- couples time
and azimuth, encoding the frame-dragging effect that is the hallmark
of a rotating spacetime.  This coupling means that no Boyer--Lindquist
observer at fixed $(r, \theta, \varphi)$ can avoid being dragged along
in the direction of the black hole's rotation; the detailed
consequences are explored in \cref{sec:kerr-horizons}.
\physnote{The cross-term $g_{t\varphi}$ is responsible for the
brightness asymmetry in images of rotating black holes: photons
co-rotating with the hole are blueshifted, while counter-rotating
photons are redshifted.}

\paragraph{Coordinate singularities.}

Two coordinate pathologies appear in the
line element~\eqref{eq:kerr-bl-metric}.  The coefficient $g_{rr} =
\Sigma/\Delta$ diverges when $\Delta = 0$, i.e.\ at the horizons
$r = r_\pm$.  This is a coordinate artefact of the same kind as the
$r = 2M$ singularity in Schwarzschild coordinates
(\cref{sec:schw-coordinates}): the Kretschmann scalar remains finite
at both horizons, and horizon-penetrating coordinates
(\cref{sec:kerr-ks}) remove the divergence.  The second pathology is
the usual axis singularity at $\sin\theta = 0$, shared with all
spherical coordinate systems.  Neither singularity is physical, but
both make Boyer--Lindquist coordinates unsuitable for numerical
integration near the horizons --- the codebase uses Kerr--Schild
coordinates instead (\cref{sec:kerr-ks}).
\warnnote{The genuine curvature singularity of the Kerr spacetime is
at $\Sigma = 0$ ($r = 0$, $\theta = \pi/2$): a ring of
coordinate radius zero in the equatorial plane.  Unlike the
Schwarzschild point singularity at $r = 0$, the Kerr singularity has
the topology of a circle.}

\paragraph{Limiting cases.}

The Boyer--Lindquist metric interpolates between two clean limits.
Setting $a = 0$ reduces $\Delta$ to $r(r - 2M)$, $\Sigma$ to $r^2$,
and $A$ to $r^4$; the line element collapses to the Schwarzschild
metric~\eqref{eq:schwarzschild-metric}.  Of the two horizon radii,
$r_+ = 2M$ becomes the Schwarzschild event horizon; $r_- = 0$
coincides with the curvature singularity and does not represent a
Cauchy horizon.

At the opposite extreme, $a = M$ gives the extremal Kerr black hole:
$\Delta = (r - M)^2$, so the two horizons merge at $r_+ = r_- = M$
and the surface gravity vanishes.  X-ray reflection spectroscopy of
accreting black holes suggests that spin parameters approaching this
limit occur in nature, with measured upper bounds near
$a/M \approx 0.998$ \cite{Thorne1974}.  For a
$10\,M_\odot$ black hole with near-maximal spin, the outer horizon
lies at $r_+ \approx M \approx 15\;\text{km}$ --- half the
Schwarzschild value.

\paragraph{Codebase representation.}

The ray tracer works in Cartesian Kerr--Schild coordinates, not
Boyer--Lindquist, so the BL radius $r$ is not a coordinate but a
derived quantity.  Given Cartesian coordinates $(x, y, z)$, the
relation between BL and Cartesian radii is
%
\begin{equation}
  r^4 - (x^2 + y^2 + z^2 - a^2)\,r^2 - a^2 z^2 = 0 \,,
  \label{eq:bl-quartic}
\end{equation}
%
a quartic in $r$ (biquadratic in $r^2$); its derivation and
properties are the subject of \cref{sec:kerr-bl-quartic}.  The
closed-form solution selects the larger of the two real roots of the
associated quadratic in $r^2$:
%
\begin{lstlisting}[language=Haskell]
blRadius :: Double -> Vec4 'Contravariant -> Double
blRadius a (Vec4 _ x y z)
  | a == 0    = sqrt (x*x + y*y + z*z)
  | otherwise =
      let p2   = x*x + y*y + z*z
          a2   = a*a
          diff = p2 - a2
          disc = diff*diff + 4*a2*z*z
          r2   = (diff + sqrt disc) / 2
      in sqrt r2
\end{lstlisting}
%
\coderef{Spacetime.Kerr}{blRadius}
The fast path for $a = 0$ recovers the Euclidean radius, confirming
the Schwarzschild limit at the code level.
\implnote{The full quartic (\cref{eq:bl-quartic}) is solved by
treating it as a quadratic in $r^2$.  The discriminant
$(p^2 - a^2)^2 + 4a^2 z^2$ is manifestly non-negative, so the square
root is always real.}
